\documentclass[12pt]{article}
\usepackage[utf8]{inputenc}
\usepackage[T1]{fontenc}
\usepackage{hyperref}
\title{Ukrainian Sign Language Language Resources for K-12 Curriculum Development}
\author{Compiled by Claude Sonnet 5 (claude-sonnet-5) for Andr\'as Kornai}
\date{}
\begin{document}
\maketitle
\section{Introduction}
Ukrainian Sign Language (USL, ISO 639-3 ukl) is the deaf community sign language of Ukraine. According to Wikipedia it is overseen by the Department of Sign Language Education at the Institute of Special Education of the National Academy of Educational Sciences (NAES) of Ukraine. Its use in educating deaf people was reintroduced only in 2006, and many schools still emphasise oral Ukrainian. Most of the materials below are PDFs in Ukrainian.
\section{Primary Educational Resources}
\begin{itemize}
\item Model curriculum for USL, grades 5--6, special secondary schools for deaf and hard-of-hearing children (2023), NAES digital library: \url{https://lib.iitta.gov.ua/id/eprint/748565}
\item Grade 5 USL textbook by Bilanova, Ivashchenko and Diakiv (2024, two parts), NAES digital library: \url{https://lib.iitta.gov.ua/id/eprint/748566}
\item A second item in the same textbook series, NAES digital library: \url{https://lib.iitta.gov.ua/id/eprint/748558}
\end{itemize}
\section{Linguistic and Reference Materials}
\begin{itemize}
\item Signbook of Ukrainian Sign Language (2021), by Kulbida, Adamiuk, Drobot and Zamsha: \url{https://lib.iitta.gov.ua/732056/}
\item Ukrainian Sign Language Dictionary for Parents (2012 article), a thematic manual for hearing parents of deaf children: \url{https://lib.iitta.gov.ua/1987}
\item Wikipedia, Ukrainian Sign Language: \url{https://en.wikipedia.org/wiki/Ukrainian_Sign_Language}
\end{itemize}
\section{Religious and Community Resources}
None were found in search results.
\section{Digital and Audio Resources}
No searchable online USL dictionary or video corpus was found in these results.
\section{Curriculum Development Context}
Materials exist for the middle grades of special schools for deaf children, but the search found none for primary grades or for hearing learners. The demographic figures on Wikipedia date from 2008--2015 and may be out of date.
\end{document}